\section{过程案例：从搬运规划到局部搜索}\label{sec:case}
\subsection{完整研究轨迹}
图 \ref{fig:trajectory} 展示四种路由在开发集上的保留最优分数。每一个更新点都对应一次候选验证完成；后续方案退步时，保留曲线维持原值。最终得分来自冻结候选后的另一组测试。

\reportfigure{routing_trajectory.pdf}{开发阶段保留最优分数的变化。纵轴相对开发集初始分数 44120 归一化，横轴为自研究启动起的实际秒数；开发集包含五个用例。}{fig:trajectory}

默认路由在约 90 秒时出现明显改进，动作 Decisions 组在接近时刻进入相同量级，随后继续提高。模型 Decisions 组的首次明显提升出现在更晚阶段；双 Decisions 组得到一次小幅改进后结束。这些曲线将最终结果展开为“何时得到改进、之后如何继续”的过程。

\subsection{动作 Decisions 组的七次候选}
本节选取动作 Decisions + 规则模型路由的一轮，进一步查看代码修改、验证结果和状态传递。c1、c2 由 Qwen 生成，代码哈希完全相同，开发分数均为 44120。两次候选尝试后，规则模型路由升级到 Sol，剩余研究沿已有候选继续。

\reportfigure{action_case.pdf}{七个候选的实际开发得分与当时保留最佳。c5 的得分下降后，选择器继续保留 c4；c6 从 c4 派生，恢复了有效搜索方向。}{fig:case}

\begin{smalltable}
\begin{tabularx}{\linewidth}{@{}lrrL@{}}\toprule
候选 & 完成验证（秒） & 实际开发分 & 算法变化\\\midrule
c1 / c2 & 12.851 / 22.792 & 44120 & Qwen 两次提交相同实现，无增益\\
c3 & 94.286 & 46506 & 依据目标堆最小箱号与预计重复搬运成本进行分段规划\\
c4 & 178.438 & 46798 & 以 c3 为起点，对不同取箱阶段做局部参数搜索\\
c5 & 244.796 & 46725 & 扩大目标堆搜索范围，得分下降\\
c6 & 311.463 & 46800 & 回到 c4，改进查询与动态规划执行，增加有效搜索次数\\
c7 & 357.026 & 46849 & 识别搜索中重新汇合的状态，复用相同后缀方案\\\bottomrule
\end{tabularx}
\end{smalltable}

\subsection{算法变化与验证信号如何衔接}
搬箱任务的代价来自“为了取出当前目标而搬开上方箱子”。直接把阻挡部分搬到另一个堆，可能很快又挡住后续目标。c3 把当前搬运与未来再次搬运联系起来，通过目标堆中箱号的信息估计后续冲突，并对可移动片段做规划。其开发分数从 44120 提高到 46506，是这一轮最大的单次变化。

c4 在已有规划基础上增加局部搜索，进一步提高到 46798。c5 扩大搜索范围后，分数反而降到 46725；评测结果将这种退步传回研究状态。后续 c6 明确选择 c4 作为父候选，转向提高单位时间内有效搜索的密度。c7 再利用状态汇合后的公共后缀减少重复计算，得到本轮最高开发分。

这条链体现了两种互补的反馈作用。裁判分数决定哪个程序进入保留状态；候选之间的结构差异则帮助后续生成定位修改方向。选择器始终保留已验证方案，使一次搜索回退能够成为后续研究的信息。

研究在剩余约 3 秒时触及 15 秒的最小请求窗口，冻结 c7。随后 15 例最终测试全部通过，总分为 139972。这条案例同时包含生成模型升级和多次算法改写，展示了整个研究过程中的状态传递。
\source{[S3] \texttt{run-20261114T085616-894b} 的候选代码、父候选记录、steps 与最终评测。}
